\chapter{FUTURE ENHANCEMENTS}

Many options exist for future improvements that will enhance the accuracy, speed, and usefulness of the BCI system. This is because the current system operates as a mere prototype with lots of limitations in terms of the hardware used and the way the signals have been processed. Improvements could be made by enhancing the circuitry, having more input channels, improved classification techniques, and realistic typing interface.

\section{Hardware Improvements}

The first possible improvement in the future would be to redesign the circuit in a proper printed circuit board instead of using a matrix board. It would decrease noise and improve connections, increasing the reliability of the system. Moreover, better shielding and grounding can be used to minimize the interference caused by the power lines and other electronic devices. The DRL circuit can also be improved with adjustable gain using a potentiometer in order to adjust noise cancellation for proper testing. Better electrodes and electrode placement will also help to collect cleaner signals.

Secondly, an increase in number of EEG channels will be a major improvement in the future. Currently the system uses only one channel to detect blinks and focus at the same time. Therefore, using more channels will provide more data about the brain and its functions, allowing to classify signals more accurately. Moreover, more channels will also allow to compare the signals across all the electrodes rather than focusing on one path, reducing confusion between blink, focus and idle states. Consequently, more complicated commands can be developed for BCI systems in the future.

\section{Signal Processing Improvements}

The signal processing stage is also open to further improvements in several aspects. The use of superior methods for digital filtering could help eliminate noise and artifacts more efficiently. Given that the current device operates within a restricted frequency band, future implementations could extend their research into other EEG bands and examine their utility in the focus determination process. In addition, better artifact removal techniques could be employed to improve the ability to deal with ocular movement noise, muscle activity and power-line interference.

A great idea for an improvement would be to perform some signal processing operations, including FFT and feature extraction, on a more powerful platform than Arduino. Considering that the use of such operations with the current processor has some limitations, using a more efficient processor could help achieve better results in analyzing the brain signals. Better sampling qualities and resolution could lead to increased accuracy in detecting small changes in the EEG signal.

\section{Classification and Machine Learning Improvements}

The present system seems to employ a relatively simple approach to classify blink, focus and idle states, however in the future more sophisticated approaches can be employed. For instance, such features of EEG signal as amplitude, band power, energy, variance or time frequency representation can be extracted and provided as input data for classification algorithms such as support vector machines, k-nearest neighbors, decision trees or neural networks.

Machine learning can also allow the system to adapt to specific user. Indeed, as EEG signals vary from person to person, a model developed for one particular user cannot necessarily be used for another person. In the future, the system can incorporate user calibration and training to adapt to a particular user. Also, the system can use adaptive threshold or even online learning to ensure more consistent performance across different sessions.

\section{Interface and functionality improvements}

Even the typing interface itself can be enhanced in further research. At present, only a basic BCI can be created as the system recognizes only a few signals. However, with improved classification and additional channels, extra functions may be introduced, such as select, delete, move left, move right, confirm or switching between different keyboard groups. These additional functions will provide a user with greater flexibility in terms of typing. In addition, predictive text or word suggestion can be applied to decrease the number of necessary selections while typing.

The user experience can also be enhanced through improvement in the speed and intuitiveness of the interface. Larger visual feedback, improved cursor control and clear indication of the current status of the signals would assist the user in recognizing the actions performed. A calibration page might also be introduced in order to adapt the system before the beginning of typing.

\section{Testing and Research Improvements}

It would also be good for future work to have more testing done. The present findings seem to be the product of a few trials. Thus, in the next iteration of the project, it would be good to conduct more tests on more trials and with more participants. That way, the reliability of the results may be established. At the same time, this will improve the validity of performance metrics obtained and comparisons that could be drawn between different versions of the software.

Finally, in future research, it would also be beneficial to test various filters, positions of electrodes, feature extractions, and classifiers. By testing such parameters, one may obtain information on how to get better performance. In particular, another version of the project might involve comparing single and multichannel BCIs. Thus, the improvements will help transform the present prototype into a more sophisticated and effective EEG-based brain-computer interface.
